% Transcription — fonds Alexandre Grothendieck, shelfmark 85, batch 2 (pages 21-40).
% Established from the facsimile archives/batches/85/batch-02.pdf, cut from a
% copy of 85.pdf fetched through the project's own relay
% (https://grothendieck-relay.onrender.com/source/85.pdf); PDF page k =
% archivists' page 20+k, no cover sheet. Per the skill transcribe-grothendieck.
% The folder is 97 pages in five batches (1-20, 21-40, 41-60, 61-80, 81-97),
% transcribed in parallel; this pass read none of the other batches.
%
% Pass: Opus 5.5 (claude-opus-5-5), 2026-09-26 — first pass, unchecked against the pages by a human.
%
% Dating. The folder is « s.d. » in the catalogue; it belongs to the
% inventory group « Géométrie et topologie combinatoire », whose range is
% given in \dating as the group's, per the skill's group rule. No date in his
% hand appears on these twenty pages.
%
% Pages skipped: none apart from two covers. Pages 22 and 32 are tan sheets
% carrying only a title in his hand, top right, and a pencilled page count
% top left: p. 22 « Invariants [underlined] des sections des formes de GP(1)
% etc » with « (9 p) », p. 32 « Sous-groupes finis de GP(1) » with « (8 p) ».
% By the cover rule they take no \page marker; their words become section
% headings, with a \note. Pages summarised: none. Everything else (21, 23-31,
% 33-40) is transcribed from his hand, in blue ink.
%
% His pagination: none on the leaves. The only numbered runs are the counts
% on the covers: « (9 p) » (pp. 23-31, exactly nine leaves) and « (8 p) »
% (pp. 33-40, eight leaves, i.e. the second run ends with this batch; p. 40
% ends mid-statement and may continue on p. 41 in batch 3). Within the runs
% he numbers cases 1)-6), a)-e), 1°)-4°) and corollaries Cor. 1-3, « Cor.
% Main » (three times on p. 34). Page 21 continues the end of batch 1 (a
% « Cor » on injective homs D_i x Z/2 -> GP(1,k), using alpha, rho, sigma,
% phi introduced there).
%
% What the batch is about. Despite the folder's title, these pages are not
% combinatorics of maps directly but the group theory behind regular maps.
% P. 21: when D_i x Z/2 embeds in GP(1,k) over an algebraically closed
% field. Pp. 23-31 (« Invariants d'automorphismes d'une droite projective »):
% for G a form of GP(1) or SL(2) over a scheme S, the invariant
% A : G -> E^1_S (the representable functor of conjugation-invariant
% functions), characterised by A(i(lambda)) = lambda + 1/lambda on a torus;
% A(1) = 2, A = -2 on involutions (« réflexions »), A = 2 on unipotents;
% conjugacy of u, u' is detected by A when A(u) +- 2 is invertible, with an
% infinitesimal counterexample at A = +-2 (p. 26); polynomials G_n with
% A(u^n) = G_n(A(u)); A(u_C) = (Tr u)^2/det u - 2 for u in GL(V), and for
% pseudo-reflections; u^n = id and order exactly n iff F_n(alpha) = 0 resp.
% Psi_n(alpha) = 0 (the « half-cyclotomic » polynomials), with corollaries
% alpha = -1, +1, 0 and alpha^2 + alpha - 1 = 0 for orders 3, 6, 4, 5; orders
% of elements of SL(2,k); a Theorem sorting n (A: n not p; B: n = 2p;
% C: n = p) by residue characteristics, flagged « à vérifier » and « peut-être
% faux ». Pp. 33-40 (« Sous-groupes finis de GP(1,k), SL(2,k), GL(2,k) »):
% in characteristic 0, finite subgroups of GP(1) are classified up to
% conjugacy by isomorphism type (Z/n, D_n, A_4, S_4, A_5 — « exactement les
% types des groupes d'isom. des cartes orientées régulières sphériques »);
% normalisers and centralisers case by case; « Cor. Main »: pairs (G, Gamma)
% over S of char. 0 correspond to regular polygons (D_n), S_4-torsors,
% tetrahedra and octahedra (D_2, A_4, S_4); models over Q; for A_5 a
% theorem: pairs (G, phi : A_5 -> G) are rigid and represented by
% Spec Z[T]/(T^2+T-1) via alpha(phi) = A(phi(pi)), pi a 5-cycle. P. 38-39:
% « Formes de D_n », dihedral extensions as torsors under abelian groups;
% Aut(D_n) = Aff(1, Z/n).
%
% Notation fixed once for the batch: GP(1) as \mathrm{GP}(1); his large
% gothic A and S (alternating and symmetric groups) as \mathfrak{A},
% \mathfrak{S}; the double-struck D_n as \mathbb{D}_n; his cursive N and Z
% for normaliser and centraliser as \mathcal{N}, \mathfrak{Z}; the affine
% line E^1_S as \mathbb{E}^1_S; the semi-direct product, which he writes as
% a dot with ½ beneath, as \cdot_{1/2}; ₂Γ (elements of order dividing 2)
% as {}_2\Gamma. « ½ simple » (p. 23) is kept as written: semi-simple.
%
% Hardest pages: 21 (a crowded corollary with overwritten cases), 26 (the
% infinitesimal computation with overwritten signs and an oblique margin),
% 29 (a slanted margin note and a cancelled block), 33 (pencil diagonals
% across the lower half), 35-36 (oblique margins, one scribbled over) and
% 40 (a boxed cancelled block). Readings to check first: p. 21 the sentence
% after « mu_2 »; p. 25 « rev. étale de rang 2 »; p. 26 « Z/p^nZ »; p. 29
% « F_n(alpha) = 0 » in 1°); p. 34 the first « Cor. Main » (ending
% « \ill{} de rg 2 sur S »); p. 39 « gerbe » (twice, above a struck
% « cat. »).


\documentclass[11pt,a4paper]{article}
\input{../preamble/grothendieck}

\folder{85}
\batch{2}
\pages{21}{40}
\dating{s.d. — le groupe « Géométrie et topologie combinatoire » (69 à 102) est daté 1976-[vers 1986]}
\watermark{Édition de démonstration}
\foldertitle{2-polyèdres et 1-polygônes : notes manuscrites (s.d.)}

\begin{document}

\page{21}
\textbf{\emph{Cor}} Sur le corps alg.\ clos $k$, considérons la condition
($i$ entier \emph{pair} $\geq 0$)
\[
  (**) \qquad \exists \text{ hom injectif } \Gamma = \mathbb{D}_i \times \mathbb{Z}/2 \longrightarrow \mathrm{GP}(1,k)
\]
\note{la parenthèse « $i$ entier pair $\geq 0$ » est entourée d'une boucle
rattachée à la fin de la première ligne. Les lettres $\alpha$, $\rho$,
$\sigma$, $\varphi$ viennent de la page précédente (lot 1)}

\emph{Alors} \add{dans $i \geq 2$,}
\begin{enumerate}
  \item[a)] Si $i =$ \struck{\ill{}} $2$, la condition \struck{(**) est
  vérifiée} \add{(**) est} vérifiée sss car.\ $k = 2$
  \item[b)] Si $i = 0$, elle est vérifiée sss ou bien car.\ $k = 0$, ou bien
  car.\ $k > 0$ et $k$ non alg.\ sur le corps premier.

  \struck{Si $i = 2$, la condition (**) est vérifiée sss car.\ $k = 2$.}
  \add{\struck{donc $\Gamma \simeq (\mathbb{Z}/2)^3$}}
  \item[c)] Supposons $i \neq 0, 2$ i.e.\ $i \geq 4$ \add{$= 2\nu$}. La
  condition pour qu'il existe mono $\mathbb{D}_i \hookrightarrow
  \mathrm{GP}(1,k)$ est que $p \nmid i$ ou \struck{$p \mid$} $i = p$
  $(\neq 2)$, dans le premier cas on a $\alpha \neq 2$ i.e.\ $\rho \neq 0$,
  donc $\mathrm{Aut}(G, \varphi|\mathbb{D}_i) \simeq \mu_2$ \add{$=$
  \struck{\ill{}}} dont l'image dans le groupe des \add{$\sigma$}
  \uncertain{facteurs} $\mathbb{Z}/2$ de $\Gamma$ est \ill{} 1,
  \struck{on} soit $u^{\nu}$. D'où le
  résultat. Dans le 2\textsuperscript{e} cas on a $\alpha = 2$, i.e.\
  $\rho = 0$, et en car.\ $\neq 2$ on a
  $\mathrm{Aut}(G, \varphi|\mathbb{D}_i) = (p)$, (donc un hom
  $\mathbb{D}_i \times \mathbb{Z}/2$ \uncertain{injectif} sur $\sigma$) et en
  car.\ 2
\end{enumerate}
\marginal{distinguer cas $i = 2$ \ill{} \ill{} \ill{} dans \ill{}}
\note{la note marginale est en face de a), marquée d'un trait vertical. Le
petit $\sigma$ est ajouté au-dessus de « des ». Le
bloc « donc $\Gamma \simeq (\mathbb{Z}/2)^3$ » est encadré, rattaché à la
ligne biffée qui répète a), elle-même couverte d'un griffonnage. Au-dessus
de « $\mu_2$ », un « $=$ » suivi d'un signe biffé ; la phrase qui suit est
lue avec peine. La page s'arrête sur « et en car. 2 » et un trait
horizontal}

\section*{\emph{Invariants} des sections des formes de $\mathrm{GP}(1)$ etc}
\note{inscrit en haut à droite d'un feuillet de garde (page 22), qui ne
porte rien d'autre ; « Invariants » est souligné. En haut à gauche, au
crayon, « (9 p) », qui compte les neuf feuillets suivants (pages 23 à 31).
Le feuillet ne reçoit pas de numéro de page}

\page{23}
\subsection*{Invariants \struck{des} d'automorphismes \struck{de} d'une droite projective}

1) Soit $G$ une forme de $\mathrm{GP}(1)$, ou de $\mathrm{SL}(2)$
\add{sur $S$} --- i.e.\ un schéma en groupes ½ simple de rang 1 sur un
schéma $S$. Considérons le foncteur
\[
  \mathbf{I}(G) = \underline{\mathrm{Hom}}_{\mathrm{sch}}(G, \mathbb{E}^1_S)^{\circlearrowleft}
  \overset{\mathrm{d\acute{e}f}}{=}
  \underline{\mathrm{Hom}}_{\mathrm{sch}}(G, \mathbb{E}^1_S)^{G}
\]
où $G$ opère sur le $\underline{\mathrm{Hom}}$ \uncertain{via} autom.\ int.\
sur $G$. Ce foncteur est représentable \add{(vrai pour $H$ \ill{} ½ simple)},
et on a un isomorphisme canonique
\[
  \mathbf{I}(G) \xrightarrow{\ A\ } \mathbb{E}^1_S
\]
qui, en tant que morphisme, peut être caractérisé par la propriété
suivante : si $\mathbb{G}_m \xrightarrow{i} G$ est un monomorphisme, alors
\[
  A(i(\lambda)) = \lambda + \lambda^{-1}
\]
et itou après tout chgt de base.
\note{l'exposant du premier $\underline{\mathrm{Hom}}$ est une petite flèche
recourbée, rendue ici par $\circlearrowleft$ ; « ½ simple » : semi-simple.
$\mathbb{E}^1_S$ est la droite affine sur $S$}

\struck{\textbf{Proposition 2} Soient $u, v$ deux sections de $G$, telles que}
\add{\struck{2) Prouvons que}} \add{on suppose}

On a donc
\[
  A(1) = 2 .
\]

[\struck{Supposons $G$ forme de $\mathrm{GP}(1)$, et} Soit $\sigma$ une
section de $G$ sur $S$ telle que $\sigma^2 = \mathrm{id}$,
$\sigma_s \neq \mathrm{id}_s$ $\forall s \in S$ \struck{(donc on
pourrait appeler $\sigma$ une réflexion de $G$)} alors
\[
  A(\sigma) = -2 .
\]
\marginal{\textbf{NB} Dans le cas d'une forme de $\mathrm{GP}(1)$, $\sigma$
est une « réflexion » ; dans le cas de $\mathrm{SL}(2)$, on \ill{} \ill{},
\ill{} générateur du \ill{}\ldots}
C'est clair si 2 inv.\ sur $S$, car alors $\sigma$ est contenu
\add{(\uncertain{marqué} dans le cas de $\mathrm{GP}(1)$ --- dans le cas de
$\mathrm{SL}(2)$, $\sigma$ est central)} dans un \struck{(unique)} tore
maximal, qu'on peut supposer scindé (par localisation), et on a alors
$\lambda = -1$, d'où $A(\sigma) = -1 + 1/{-1} = -2$. Pour le cas général, on
fait un argument de coeff universels\ldots]

\struck{Prop} \struck{\uncertain{Remarque}} Supposons que $u$ soit
\emph{unipotent}, par quoi on entend qu'il existe \struck{loc.}
\uncertain{fpqc} une \struck{forme}

\page{24}
monomorphisme
\[
  \mathrm{Aff}(1)_S = \mathbb{G}_m \cdot_{1/2} \mathbb{G}_a \xrightarrow{\ i\ } G
\]
et une section $t$ de $\mathbb{G}_a$, telle que $u = i(t)$.
\note{le « ½ » est écrit sous le point du produit : produit semi-direct}
Je dis qu'alors aussi on a
\[
  A(u) = 2
\]
En effet, on trouve plus \struck{précisément} \add{généralement}
\[
  A(i(\lambda \cdot t)) = \lambda + \lambda^{-1}
\]
On le voit p.\ ex.\ en notant que (pour $\lambda$, $t$ indét.)
\[
  f(\lambda, t) \overset{\mathrm{df}}{=} A(i(\lambda, t))
\]
étant \uncertain{élém.} $\in \underline{O}_S[\lambda, t][\lambda^{-1}]$
satisfaisant les conditions
\[
  \begin{cases}
    f(\lambda, 0) = \lambda + \lambda^{-1} \\
    f\bigl(\underbrace{(0,t)(\lambda,0)(0,-t)}_{(\lambda,0)(\lambda^{-1}t - t) = (\lambda, \lambda^{-1}t - t)}\bigr) = \lambda + \lambda^{-1}
  \end{cases}
\]
\note{« plus gén. » est écrit devant la seconde ligne de l'accolade}
donc, posant $s = \lambda^{-1}t - t$ i.e.\ $t = \dfrac{s}{\lambda^{-1} - 1}$,
on a
\[
  f(\lambda, s) = \lambda + \lambda^{-1}
\]
dès que $\lambda^{-1} - 1$ inversible --- donc par prolongement, pour tt
système $(\lambda, s)$.

Une condition \emph{nécessaire} pour que $u$ soit unipotent est donc que
$A(u) = 2$. On voit aisément que cette condition est suffisante au-dessus
d'un corps (utilisant le fait que sur un corps alg.\ clos, tt élément de $G$
est contenu dans un \emph{s-groupe} de Borel). J'ignore ce qui en est dans le
cas général.
\note{le « 2 » de « $A(u) = 2$ », étroit et appuyé, pourrait se lire 1}

\page{25}
\textbf{\emph{Proposition}} Soient $u, u'$ deux sections de $G$,
\struck{telles} \add{on} suppose $A(u) - 2$ \add{et $A(u) + 2$} inv.\
[i.e.\ $u_s^2$ non unipotent $\forall s \in S$]. Pour que $u$ et $u'$
soient conjugués loc.\ \uncertain{fpqc} (ou encore, conjugués si $G$ forme
de $\mathrm{GP}(1)$) il f.\ et s.\ que $A(u) = A(u')$.
\add{Ceci vaut aussi sans restriction sur $A(u)$, si $S$ spectre d'un corps
\uncertain{et $u_s$ non central $K_S$}}
\note{le paragraphe de l'énoncé est marqué d'un trait vertical en marge
gauche}

\emph{Néc.}\ évident, prouvons suff.\ sur $u$ et $u'$. \struck{En effet},
l'hypothèse implique que $u$, $u'$ sont contenus chacun dans un (unique)
tore maximal. Utilisant la th.\ de conjugaison, on peut supposer que c'est
le même, et par localisation, que c'est $\mathbb{G}_m$. Donc $u$, $u'$
correspondent à des sections $\lambda$, $\lambda'$ de $\mathbb{G}_m$, et
l'hyp $A(u) = A(u')$ signifie
\[
  (*) \qquad \lambda + \lambda^{-1} = \lambda' + \lambda'^{-1} = \alpha \in \Gamma(S, \underline{O}_S)
\]
ce qui signifie que $\lambda$ et $\lambda'$ sont \struck{\ill{}} solutions
de la même équation
\[
  (**) \qquad t^2 - \alpha t + 1 = 0
\]
dont le discriminant est
\[
  \alpha^2 - 4 = \underbrace{(\alpha - 2)}_{\text{inv.}}(\alpha + 2)
\]
\struck{il résulte de ceci que aux pts où $\alpha(s) = -2$. L.} dont inv.\
par hyp.\ \struck{$\alpha^2 - 4$ inv.\ i.e.\ $\alpha + 2$ inv.\ i.e.\
$\alpha(s) \neq -2$ en tt $s \in S$, donc} Cette équation définit un
\uncertain{rev.\ étale de rang 2} de $S$, et \struck{on trouve que l'on a
loc.\ deux solutions} deux sections \struck{distinctes} \ill{} sont égales
loc., soit conjuguées par l'involution canonique $t \mapsto t^{-1}$, on gagne
donc. Il reste à faire une étude locale \struck{\ill{}} plus serrée aux pts
où $\alpha(s) = \pm 2$, i.e.\ $\lambda_s = \lambda'_s = \pm 1$, prenons le cas
$-1$ p.ex.

\page{26}
Les réductions habituelles nous ramènent au cas d'un anneau \add{de base}
\struck{artinien \ill{}, corps} \uncertain{$\mathbb{Z}/p^n\mathbb{Z}$}
artinien $A$, corps résiduel \struck{$k$} $\mathbb{Z}/p\mathbb{Z}$, i.e.\
prouvons que si $\lambda \neq \lambda' \in A^{*}$ sont tels que
$\lambda + \lambda^{-1} = \lambda' + \lambda'^{-1}$, alors on a
$\lambda' = \lambda$ ou $\lambda' = \lambda^{-1}$. Supposons qu'on sache
$\lambda'_n = \lambda_n$, prouvons $\lambda'_{n+1} = \lambda_{n+1}^{\pm 1}$ ?
\note{sous $\lambda'_{n+1}$ et $\lambda_{n+1}$, des signes « $=$ » verticaux
renvoient à $\lambda'$ et $\lambda$. Le « $\neq$ » de « $\lambda \neq
\lambda'$ » est peut-être un « $=$ » surchargé ; les lignes « artinien »,
« $\mathbb{Z}/p^n\mathbb{Z}$ » et « $k$ » sont enchevêtrées de ratures}

On a \struck{$\lambda'_n = \lambda_n$} $\lambda' = \lambda + \mu$, avec
$\mathfrak{m}\mu = 0$, \struck{d'où} donc
\begin{align*}
  \lambda' &= \lambda(1 + \lambda^{-1}\mu) = \lambda(1 - \mu) = \lambda + \mu \\
  \lambda'^{-1} &= \lambda^{-1}\underbrace{(1 - \mu)^{-1}}_{1 + \mu} = \lambda^{-1}(1 + \mu) = \lambda^{-1} - \mu \\
  \lambda' + \lambda'^{-1} &= \lambda + \lambda^{-1} + \mu(\lambda^{-1} - \lambda) = \lambda + \lambda^{-1} .
\end{align*}
\note{les signes devant le dernier $\mu$ de la deuxième ligne et devant
$\mu(\ldots)$ à la troisième sont surchargés (un signe noirci sur un
autre) ; l'exposant du second $\lambda$ de la parenthèse est surchargé,
illisible}
\marginal{\textbf{NB} Il s'agit \ill{} d'une propriété générale des
morphismes \struck{\ill{}} « ramifiés » (non étales) en \ill{} : \ill{}
« \ill{} » des \ill{} des solutions infinitésimales \ill{}\ldots)}
Donc l'hyp $\lambda' + \lambda'^{-1} = \lambda + \lambda^{-1}$ résulte de
$\lambda'_n = \lambda_n$ i.e.\ $\lambda' = \lambda + \mu$,
$\mathfrak{m}\mu = 0$ ; \ill{} \ill{} \ill{} \uncertain{autant}
\[
  \lambda' = \lambda \quad \text{ou} \quad \lambda'^{-1} = \lambda \ ?
\]
\struck{i.e.} \struck{On} ces relations signifient, \struck{\ill{}}
\add{l'une} $\mu = 0$, l'autre \struck{$\lambda^{-1}(1+\mu) = \lambda$ i.e.\
$\lambda\mu = 0$} \struck{$\lambda^{-1}(1+\mu)$ \quad
$\lambda - \lambda^{-1} + \mu = 0$}
\[
  \mu = \lambda^{-1} - \lambda
\]
donc ne correspondent qu'à deux valeurs possibles \add{au plus} de $\mu$ ---
donc (si $k$ n'est pas réduit à 2 él.)\ on trouve des solutions $\mu$ de
façon à mettre en défaut la relation
$[(\lambda' = \lambda) \text{ ou } (\lambda' = \lambda^{-1})]$.

\page{27}
\textbf{\emph{Remarque}} Soit $j : \mathbb{G}_m \to G$ un hom.\ de groupes,
pas nécessairement \struck{injectif} \add{mono}, supposons que le noyau soit
$\mu_n$, donc $j$ se factorise en

\begin{tikzcd}
  \mathbb{G}_m \arrow[r, "\lambda \mapsto \lambda^n"] \arrow[rr, bend right=30, "j"'] & \mathbb{G}_m \arrow[r, "i"] & G
\end{tikzcd}

On en conclut que
\[
  A(j(\lambda)) = \lambda^n + \lambda^{-n} = G_n(\lambda + \lambda^{-1})
\]
où
\[
  G_n \in \mathbb{Z}[U]
\]
est un polynôme unitaire universel de degré $n$, [celui exprimant
$T^n + T^{-n} \in \mathbb{Z}[T, T^{-1}]^{\sigma}$ $(\sigma T = T^{-1})$
$\simeq$ \struck{\ill{}} $\mathbb{Z}[U]$, $U = T + T^{-1}$]
\note{l'isomorphisme $\mathbb{Z}[T,T^{-1}]^{\sigma} \simeq \mathbb{Z}[U]$
est écrit verticalement}

\marginal{\uncertain{remarque}}
\emph{Exemples} 1) Soit
$\mathbb{G}_m \overset{i}{\hookrightarrow} \widetilde{G}$ (forme de
$\mathrm{SL}(2)$) $\xrightarrow{\ p\ } G = \widetilde{G}/\mu_2$ (centre de
$\widetilde{G}$) on a donc
\begin{align*}
  A(i(\lambda)) &= \lambda + \lambda^{-1} = \alpha \\
  A(p\,i(\lambda)) &= \lambda^2 + \lambda^{-2} = \alpha^2 - 2
\end{align*}
Plus généralement, indépendamment \ill{} \ill{} \ill{}
\struck{2) Soit \ill{} un pt.\ \ill{} dans $V$ (\ill{} \ill{} 2 \ill{})}
pour toute section $u$ de $\widetilde{G}$ sur $S$
\[
  \begin{cases}
    A'(p(u)) = A(u') = A(u)^2 - 2 \\
    A(u) = \mathrm{Tr}\, u
  \end{cases}
\]
[se ramène au cas précédent]

2) Soit $V$ Mod.\ loc.\ libre rg 2 sur $S$, ($u$ automorphisme de $V$,
\struck{alors on} (pas nécessairement dans $\mathrm{SL}(V)$) $u_C$
l'automorphisme de $C$ \add{\uncertain{$C = \mathbb{P}(V)$}} qu'il définit,
alors on a
\[
  A(u^C) = \frac{(\mathrm{Tr}\, u)^2}{\det u} - 2
\]
(se ramène au cas 1), en écrivant loc.\ $u = tv$, $t$ unité,
$v \in \mathrm{SL}(V)$\ldots)
\note{« $C = \mathbb{P}(V)$ » est ajouté au-dessus de la ligne ; on lirait
plutôt « $\mathbb{P}(C)$ »}

Si p.ex.\ l'équation aux v.\ propres de $u$
\[
  t^2 - (\mathrm{Tr}\, u)t + \det u = 0
\]

\page{28}
se décompose en $(t - \lambda_1)(t - \lambda_2)$, i.e.
\[
  \lambda_1 + \lambda_2 = \mathrm{Tr}\, u, \qquad \lambda_1\lambda_2 = \det u
\]
alors on trouve
\[
  A(p(u)) = \frac{\lambda_1}{\lambda_2} + \frac{\lambda_2}{\lambda_1} = \frac{\lambda_1^2 + \lambda_2^2}{\lambda_1\lambda_2}
\]
P.ex.\ si $u$ est une \emph{pseudo-réflexion}, \struck{\ill{}} de paramètre
$\lambda = \det u$ --- on a les « v.\ propres » $\lambda$ et $1$, on trouve
\[
  A(pu) = \lambda + \lambda^{-1} \qquad (u \text{ ps.\ réfl.\ par } \lambda)
\]
On peut aussi le voir en notant que une ps.\ réflexion s'insère dans une
famille à un paramètre, correspondant à
\[
  \mathbb{G}_m \longrightarrow \mathrm{GL}(V) \longrightarrow \underline{\mathrm{Aut}}_C
\]
et appliquer la définition générale de l'invariant.
\note{sous $\mathbb{G}_m \to \mathrm{GL}(V)$ une accolade, marquée
« mono »}

\textbf{\emph{Proposition}} Soit $u$ une section de $G$, \struck{Pour que}
\add{telle que} $u_s \neq \mathrm{id}$, $\forall s \in S$.
\struck{$u^2 = 1$, il \ill{}} Si $u^2 = 1$, \add{(et supposons car.\ rés.\
$\neq 2$ si $G$ forme de $\mathrm{SL}(2)$,)} on a
\[
  A(u) = -2 ,
\]
[la réciproque étant vraie si $S$ est réduit \add{\uncertain{2 inv., par}}]

\textbf{\emph{Dém}} Dans le cas \struck{de} \add{d'une forme de
$\mathrm{GP}(1)$,} \add{$u_s \neq 1_s$ \struck{$\forall s \in S$},
$u_s^2 = \mathrm{id}$ $\forall s \in S$} l'hypothèse implique que $u$
appartient à un unique tore maximal --- pour une forme de
\struck{$\mathrm{SL}(2)$ canonique de} $\mathrm{SL}(2)$, pour $u$ la section
\uncertain{gén.}\ du centre $\mu_2 = \lbrace \pm 1 \rbrace_S$, dans un tore,
on a $u$ \struck{\ill{}} \add{est} une section d'un tore maximal
$T \simeq \mathbb{G}_m \hookrightarrow G$, correspondant à la section $-1$ de
$\mathbb{G}_m$, d'où $A(u) = -2$. Dans le cas d'une forme de
$\mathrm{GP}(1)$, on

\page{29}
conclut le résultat \ill{} sans restr.\ sur les car.\ résiduelles,
\uncertain{car} le schéma des \struck{éléments} « réflexions » dans
$\mathrm{GP}(1)_{\mathbb{Z}}$ est lisse sur $\mathrm{Spec}\,\mathbb{Z}$. Dans
le cas d'une forme de $\mathrm{SL}(2)$, OPS que $G = \mathrm{SL}(V)$,
\struck{i.e.}\ \add{i.e.}\ $u$ est un autom de $V$ satisfaisant
\struck{\ill{}}, $\det u = 1$, \struck{Mais $u_s \neq \mathrm{id}$
$\forall s \in S$} Comme, d'après Hamilton--Cayley
\struck{\ill{}} $u^2 - (\mathrm{Tr}\, u)u + \mathrm{id} = 0$, donc
\add{$u^2 = \mathrm{id} \Longleftrightarrow$} $(\mathrm{Tr}\, u)u =
2\,\mathrm{id}$, \struck{\ill{}} \add{et cette hyp.} \add{il faut conclure
$\mathrm{Tr}\, u = -2$.} \struck{\ill{}} Faisant un calcul infinitésimal
\struck{\ill{} \ill{} \ill{} \ill{} \ill{}, \ill{} \ill{} \ill{} \ill{}
\ill{} \ill{} \ill{} \ill{} nécessité de cette restriction.}
\marginal{\ill{} \ill{}}
car.\ résiduelle 2, OPS que $\mathrm{Tr}\, u_n = -2$, et \add{\ill{}
\ill{}} prouvons $\mathrm{Tr}\, u_{n+1} = \mathrm{Tr}\, u = -2$.
\struck{On \ill{} $\mathrm{Tr}\, u + 2 = \tau$ \ill{} \ill{}}
\note{le passage « Faisant un calcul infinitésimal \ldots{} $= -2$ » est
marqué d'un trait vertical en marge, où deux mots ne sont pas lus}

\textbf{\emph{Proposition}} Soient $G$ une forme de $\mathrm{GP}(1)_S$ ou
$\mathrm{SL}(2)_S$, $u$ une section de $G$, \add{$\alpha = A(u)$,}
\struck{telle que $(A(u) + 2)(A(u) - 2)$ soit inv.,} $n$ un entier
$\geq 3$. \struck{Pour que l'on ait Soient $F_n$} \add{On suppose
$(\alpha + 2)(\alpha - 2)$ inv.}, Soient
\[
  F_n \in \mathbb{Z}[U], \qquad \Psi_n \in \mathbb{Z}[U]
\]
les deux polynômes \struck{qui} unitaires \add{à racines dans
$\overline{\mathbb{Q}}$ \struck{distinctes}} dont les racines sont
1°) $\zeta + \zeta^{-1}$, où $\zeta$ parcourt les racines $n$-ièmes de 1
\struck{telles} distinctes de $\pm 1$ 2°) les $\zeta + \zeta^{-1}$, où
$\zeta$ parcourt les racines \emph{primitives} $n$-ièmes de 1.

Sous ces conditions
\begin{enumerate}
  \item[1°)] Pour que $u^n = \mathrm{id}$, il f.\ et s.\ que
  \uncertain{$F_n(\alpha) = 0$}
  \item[2°)] Pour que \add{\uncertain{i.e.\ soit}} $u^n = \mathrm{id}$ et que
  $u_s$ soit d'ordre $n$ exactement pour tt $s \in S$, il f.\ et s.\ que
  $\Psi_n(\alpha) = 0$.
\end{enumerate}
\note{le dernier signe de 1°) ressemble à un « 1 » ; « Sous ces conditions »
suit un « P » souligné, relié par une flèche aux deux alinéas}
\marginal{Sans hyp.\ \ill{} \ill{} que $\alpha^2 - 4$ inv., \ill{} dire que
[$u^n = \mathrm{id}$ \struck{\ill{}} et les $u_s$ d'ordre $n$ exactement
\add{premier aux car.\ résiduelles}] $\Longleftrightarrow$
[$\Psi_n(\alpha) = 0$ \emph{et} $\alpha^2 - 4$ inv.]}

Principe de démonstration. L'hyp.\ \struck{\ill{}} \add{$\alpha \neq 2$}
inv.\ \add{i.e.\ les $u_s$ non unipotents,} implique déjà l'existence d'un
unique tore maximal $T$ qui contient $u$. OPS $T \simeq \mathbb{G}_m$, donc
$u$ donné par une

\page{30}
section $\lambda$ de $\mathbb{G}_m$, donc \struck{$*$}
\[
  \alpha = \lambda + \lambda^{-1} ,
\]
et on est ramené à une question de calcul standard.

\struck{\uncertain{Remarques}}

\textbf{\emph{Cor.\ 1}} Supposons les car.\ résiduelles de $S \neq 3$. Pour
qu'une section $u$ de $G$ satisfasse $u^3 = \mathrm{id}$ (resp.\
$u^6 = \mathrm{id}$) et que l'ordre des $u_s$ $(s \in S)$ soit toujours égal
à 3 (resp.\ 6) il f.\ et s.\ qu'on ait $\alpha = -1$ (resp.\ $\alpha = +1$).

\struck{L'une et l'autre hyp.\ implique que}

\textbf{\emph{Cor 2}} Supposons les car.\ rés.\ de $S \neq 2$. Pour que l'on
ait $u^4 = \mathrm{id}$ et que les $u_s$ $(s \in S)$ soient d'ordre 4
exactement, il f.\ et s.\ que $\alpha = 0$

\textbf{\emph{Cor.\ 3}} Supposons les car.\ résiduelles $\neq 5$. Pour que
$u^5 = \mathrm{id}$ et les $u_s$ d'ordre 5, il f.\ et s.\ qu'on ait
\[
  \alpha^2 + \alpha - 1 = 0
\]
\note{on lit « $u^5 = 0$ », pour $u^5 = \mathrm{id}$}

\textbf{NB} L'ordre \add{$\nu$} d'un élément \add{$\tilde{u}$} de
$\mathrm{SL}(2,k)$ ($k$ un corps) --- donc d'un $\widetilde{G}(k)$, pour une
forme \add{$\widetilde{G}$} de $\mathrm{SL}(2)$ sur $k$ --- \ill{} peut être
d'une des formes suivantes \struck{\ill{}} pour une car.\ donnée :
\begin{enumerate}
  \item[a)] \add{$\nu =$} $+\infty$
  \item[b)] $\nu$ entier $\geq 1$ \emph{premier} à $p$ :
  \item[c)] $\nu = 2p$ ($p$ premier \emph{impair}) \struck{\ill{}}
  \item[d)] $\nu = p$
\end{enumerate}
\marginal{si $p =$ car.\ $k$ \uncertain{fini}}
\note{l'annotation marginale, embrassant c) et d) par une accolade, est
reliée par une flèche à d)}

Si on exclut les cas $\pm\mathrm{id}$ (élément du centre, i.e.\ d'un élément
qui donne 1 dans $\widetilde{G}/\mu_2 = G$) [qui donnent lieu à
$\alpha = \pm 2$, mais d'autres cas y donnent lieu aussi] on a vu que la
valeur de $\alpha$ détermine (loc.\ fpqc) la classe de conj.\ de $u$. On
trouve en fait l'analyse suivante :

\page{31}
\marginal{On suppose $u$ \struck{\ill{}} non central ; $\zeta$ une solution de
$\zeta^2 - \alpha\zeta + 1 = 0$ \uncertain{dans $\bar{k}$}}
\note{« $u$ non central » est souligné deux fois}

\struck{$\nu$ premier à $p$ ($p$ fini) $\Longleftrightarrow$}

$k$ \emph{car.\ 0} \emph{($p = +\infty$)} :
\begin{itemize}
  \item[\struck{a)}] $\nu$ fini $\Longleftrightarrow$ $\exists$ \add{$n \geq 3$}
  tel que $\Psi_n(\alpha) =$
  \item[\struck{b)}] \struck{\ill{}} $\Longleftrightarrow$ \struck{$\zeta$
  racine} $\zeta$ d'ordre fini $\neq 1, 2$
\end{itemize}

$k$ car.\ $p > 0$ :
\begin{itemize}
  \item[b)] $\nu$ fini $\Longleftrightarrow$ $\alpha$ algébrique
  $\Longleftrightarrow$ $\exists$ $n \geq 1$ tel que $\Psi_n(\alpha) = 0$
  ($\Longleftrightarrow$ $\exists$ $n \geq 1$, premier à $p$ ou de la forme
  $2p$, \struck{$2p$ (avec} $p$ impair) ou $p$, tel que $\Psi_n(\alpha) = 0$
  [\uncertain{\emph{n} est} alors unique, et égal à $\nu(u)$] $(S \neq
  \emptyset)$ $\Longleftrightarrow$ $\zeta$ d'ordre fini (en fait).
\end{itemize}
\note{disposé en tableau : « $k$ car.\ 0 » et « $k$ car.\ $p > 0$ » à
gauche, les équivalences à droite. Le second membre de
« $\Psi_n(\alpha) = $ » est laissé en blanc ; « $(S \neq \emptyset)$ » est
écrit à gauche du crochet}

\textbf{\emph{Théorème}} $\widetilde{G}$ forme de $\mathrm{SL}(2)_S$, $u$
section de $\widetilde{G}$, $\alpha = A(u) \in \Gamma(S, \underline{O}_S)$,
$n \in \mathbb{N}$, $\boxed{n \geq 3}$. On s'intéresse à la traduction, en
termes de $\alpha$ et des propriétés des car.\ résiduelles de $S$, de la
propriété suivante
\[
  (*) \qquad u^n = \mathrm{id}, \text{ et } \forall s \in S, \text{ l'ordre de } u_s \text{ est égal à } n.
\]
\marginal{\textbf{NB} On ne prend pas $n = 2$, qui \uncertain{se calculerait}
complètement}

A) Supposons que $n$ ne soit \struck{ni de \ill{} premier, ni $2p$ ($p$
premier \ill{})} \add{de la forme $p$ ($p$ premier \ill{} \emph{impair})}.
Alors $(*)$ équivaut :
\begin{itemize}
  \item[a)] $\Psi_n(\alpha) = 0$, et $n$ inv.\ sur $S$ (i.e.\ les car.\
  résiduelles sont premières à $n$)
\end{itemize}
\marginal{\textbf{NB} Le cas $n = 4$ \ill{} \ill{} \ill{} --- A) doit
marcher \ill{} dans le cas de $G = \widetilde{G}/\mu_2$}

B) $n = 2p$ ($p$ premier impair). Alors \ill{} équivaut :
\begin{itemize}
  \item[b)] $\Psi_n(\alpha) = 0$, et 2 inv.\ sur $S$ (i.e.\ car.\ rés.\ de
  $S$ sont $\neq 2$)
\end{itemize}

C) $n = p$, $p$ est premier \add{impair}. Alors ces équivalent :
\begin{itemize}
  \item[c)] $\Psi_n(\alpha) = 0$.
\end{itemize}
\marginal{à vérifier}
\marginal{Peut-être faux (bons contre-exemples \ill{} \ill{}\ldots)}
\note{B) et C) sont embrassés à gauche par une accolade qui porte « à
vérifier » ; l'autre note marginale est reliée à C). La note sur $n = 4$
est encerclée et rattachée à A)}

\textbf{NB} Ces conditions impliquent \struck{\ill{}}
$(\alpha + 2)(\alpha - 2)$ inv., sauf dans le cas A. Dans les cas B
\struck{resp.\ C} \add{et C}, les pts où $\alpha_s^2 - 4 = 0$ sont ceux de
car $p$, \struck{(où $\alpha_s = -2 \neq 2$] (resp.\ ceux de car \ill{}}
où $\alpha_s = -2 \neq +2$ (cas B) resp $\alpha_s = +2 \neq -2$ (cas C)
\note{le dernier « $\neq -2$ » est surchargé}

\section*{Sous-groupes finis de $\mathrm{GP}(1)$}
\note{inscrit en haut à droite d'un feuillet de garde (page 32), qui ne
porte rien d'autre. En haut à gauche, au crayon, « (8 p) », qui compte les
huit feuillets suivants (pages 33 à 40, fin du lot). Le feuillet ne reçoit
pas de numéro de page}

\page{33}
\subsection*{Sous-groupes finis de $\mathrm{GP}(1,k)$, $\mathrm{SL}(2,k)$, $\mathrm{GL}(2,k)$}
\note{titre souligné, de sa main, en tête de la page 33}

On suppose d'abord $k$ corps de car.\ 0, alg clos.

a) \emph{Sous groupes finis de $\mathrm{GP}(1,k) = G$}

\textbf{\emph{Thm}} a) Pour que deux ss-groupes finis \add{$\Gamma, \Gamma'$}
\struck{de} $G$ soient conjugués dans $G$, il f.\ et s.\ qu'ils soient
isomorphes.

b) Les classes d'isom.\ de \emph{ss}-groupes finis de $G$ \add{\uncertain{allez-y}}
sont $\mathbb{Z}/n\mathbb{Z}$ ($n \geq 1$ \struck{\ill{}}), $\mathbb{D}_n$
($n \geq 2$), $\mathfrak{A}_4$, $\mathfrak{S}_4$, $\mathfrak{A}_5$, i.e.\
(ce sont \emph{exactement} les types des groupes d'isom.\ des
\marginal{cartes orientées régulières sphériques ! (comme leurs \ill{} les
\ill{} \ill{} dans $\mathrm{GP}(1)$\ldots.)}
\note{la note marginale est reliée par un trait à « i.e. » ; la parenthèse
du texte reste ouverte}

\struck{c) Si $\Gamma$ est un sous-groupe de $G$, son normalisateur}
\add{\uncertain{et si}} il est égal à son normalisateur, sauf dans le cas de
$\mathbb{Z}/n\mathbb{Z}$ ($n \geq 1$) \struck{et} de $\mathbb{D}_2$,
\add{de $\mathfrak{A}_4$.)} Dans ces cas, \ill{}, mettant à part le cas
trivial de $1 = \mathbb{Z}/n\mathbb{Z}$ \ill{} :
\begin{enumerate}
  \item[1°)] \struck{Normalisateur de} $\mathcal{N}(\mathbb{Z}/n\mathbb{Z})$
  \struck{est le normalisateur} \add{$\mathcal{N}(T)$}, où
  $T = \mathfrak{Z}(\mathbb{Z}/n\mathbb{Z})$ est l'unique tore maximal qui
  contient $\mathbb{Z}/n\mathbb{Z}$, \struck{\ill{}}
  $\mathcal{N}(\mathbb{Z}/n\mathbb{Z})/\mathfrak{Z}(\mathbb{Z}/n\mathbb{Z}) = \mathbb{Z}/2\mathbb{Z}$
  \item[2°)] \struck{Normalisateur de} $\mathcal{N}(\mathbb{D}_2) \xrightarrow{\sim} \mathfrak{S}_4$,
  \struck{Centralisateur de} $\mathfrak{Z}(\mathbb{D}_2) = \mathbb{D}_2$,
  \[
    \mathcal{N}(\mathbb{D}_2)/\mathfrak{Z}(\mathbb{D}_2) \xrightarrow{\sim} \mathrm{Aut}(\mathbb{D}_2) \simeq \mathfrak{S}_3
  \]
  \item[3°)] $\mathcal{N}(\mathfrak{A}_4) = \mathfrak{S}_4$,
  $\mathfrak{Z}(\mathfrak{A}_4) = (1)$,
  $\mathcal{N}(\mathfrak{A}_4)/\mathfrak{Z}(\mathfrak{A}_4) = \mathfrak{S}_4 = \mathrm{Aut}(\mathfrak{A}_4)$
\end{enumerate}
\struck{D'autre part, dans tous les cas sauf}
\note{$\mathcal{N}$ et $\mathfrak{Z}$ rendent ses lettres cursives pour
normalisateur et centralisateur ; $\mathfrak{A}_4$, $\mathfrak{S}_4$ ses
grandes lettres gothiques. Des traits obliques au crayon, parallèles et
croisés, couvrent la moitié inférieure de la page, du bloc c) jusqu'au
bas ; à gauche, un dessin au crayon et à l'encre (des arcs, les marques
« 0 », « 1 », « J' »), non redessiné. On ne peut dire si ces traits
biffent le bloc c) ou appartiennent à une figure ; le texte à l'encre,
lui, reste lisible et n'est pas biffé à la plume, sauf la ligne « D'autre
part \ldots{} sauf »}

c) Pour toute carte \struck{\ill{}} orientée finie sphérique (connexe)
$C$, considérons sa réalisation topologique --- on trouve un foncteur de
groupoïdes
\[
  C \longrightarrow \mathcal{C} \longrightarrow \text{catégorie des couples } (G, \Gamma)
\]
\note{sous $\mathcal{C}$ : « cat.\ des \emph{« cartes holomorphes
sphériques »} i.e.\ des « cartes holomorphes » \struck{\ill{}}
telles que la variété complexe sous-jacente soit \struck{\ill{}} de type
de genre 0 » ; la catégorie d'arrivée est celle des couples $(G, \Gamma)$
« d'une forme de $\mathrm{GP}(1, \mathbb{C})$, et d'un sous-gr.\ fini ».
Le premier $C$ est une lettre ronde, le second, plus grande, est repassé}

\page{34}
1) Soit $\Gamma \simeq \mathbb{Z}/n\mathbb{Z}$, $n \geq 2$. Alors
$T = \mathfrak{Z}(\Gamma)$ est un tore maximal de $G$, et c'est l'unique
tore maximal de $G$ contenant $\Gamma$. On a \struck{\ill{}}
$\Gamma = {}_nT$, \struck{$\mathrm{Aut}(\Gamma) =$} et les relations
$\Gamma \mapsto T = \mathfrak{Z}(\Gamma)$, $T \mapsto {}_nT = \Gamma$,
définissent des bijections inverses l'une de l'autre entre l'ens.\ des
tores maximaux de $G$, et l'ens des ss-groupes de $G$ \struck{\ill{}
$\mathbb{Z}$} cycliques d'ordre $n$. $(\Longrightarrow)$
\[
  \mathrm{Aut}(G, \Gamma) \simeq \quad
  \bigl(\mathcal{N}(\Gamma) = \mathcal{N}(T) = , \quad
  \mathcal{N}(\Gamma)/\mathfrak{Z}(\Gamma) = \mathrm{Aut}(G, \Gamma) \simeq \mathbb{Z}/2\mathbb{Z}
\]
\note{au-dessus de « $\mathcal{N}(T) =$ », un « $\mathrm{Aut}(G, T) =$ »
biffé ; au-dessus de $\mathbb{Z}/2\mathbb{Z}$, « $\mathrm{Aut}\,\Gamma$ »
avec une flèche montante d'inclusion. À gauche, en marge, deux dessins : une
étoile de huit arêtes issues d'un sommet, et plus bas un cercle marqué de
deux points ; non redessinés}

\marginal{\textbf{NB} 1° à 3° \ill{} les seuls cas \ill{} \ill{}}
3) Soit $\Gamma = \mathbb{D}_2$ \add{$\simeq \mathbb{Z}/2 \times \mathbb{Z}/2$}.
Alors
\[
  \mathfrak{Z}(\Gamma) = \Gamma, \qquad \mathcal{N}(\Gamma) = \mathfrak{S}_4 ,
\]
$\mathcal{N}(\Gamma)$ est l'\uncertain{unique} ss-groupe du type
$\mathfrak{S}_4$ qui contienne $\Gamma$, et
$\Gamma \mapsto \mathcal{N} = \mathcal{N}(\Gamma)$,
$\mathcal{N} \longmapsto \Gamma$ = unique ss-grpe.\ du type $\mathbb{D}_2$
de $\mathcal{N}$ [= unique ss-groupe invariant \struck{\ill{}} de
$\mathcal{N}$ distinct de $[\mathcal{N}, \mathcal{N}]$ i.e.\ d'indice
$\neq 2$ [ou encore d'indice 6]], sont des corr.\ 1-1 entre l'ens des
sous-groupes de $G$ de type $\mathbb{D}_2$, et l'ens des ss-groupes de type
$\mathfrak{S}_4$.
\note{le « 3) » (peut-être un « 2) » surchargé) est pris dans un crochet
auquel mène une flèche partie de la note marginale, qui se termine par un
mot non lu}

\marginal{\textbf{NB} Si $\mathcal{N}(\Gamma) \to \mathrm{Aut}(\Gamma)$ non
injectif (c'est le seul cas \ill{} les cas)}
\[
  \mathcal{N}(\Gamma)/\mathfrak{Z}(\Gamma) \xrightarrow{\sim} \mathrm{Aut}(\Gamma) \simeq \mathfrak{S}_3
\]

\textbf{\emph{Cor.\ Main}} La donnée, sur un schéma $S$ de car.\ 0, d'une
forme $G$ de $\mathrm{GP}(1)$ et d'un sous-groupe étale fini $\Gamma$ de $G$
cyclique d'ordre $n$ ($n \geq 2$ donné) équivaut à celle d'une
\struck{Mod.\ \ill{}} \add{\ill{} \ill{}} \ill{} de rg 2 sur $S$.
\marginal{\ill{} \ill{} \ill{} \ill{}}
\note{ce corollaire est écrit à droite du 3), serré dans la marge droite,
marqué d'un trait vertical ; la ligne « $\mathcal{N}(\Gamma)$ est
l'unique ss-groupe du type \ldots{} » passe à travers sa fin}

\textbf{\emph{Cor.\ Main}} \struck{La donnée d'un ss-groupe}

\textbf{\emph{Cor.\ Main}} \add{Sur base $S$ de car.\ 0} La donnée d'un
$(G, \Gamma)$, $G$ forme de $\mathrm{GP}(1)$, \struck{$\Gamma$ \ill{}}
$\Gamma \subset G$ forme de $\mathbb{D}_2$, équivaut à celle des
$(G, \widetilde{\Gamma})$, $G$ \ldots, $\widetilde{\Gamma} \subset G$ forme de
$\mathfrak{S}_4$, ou encore à celle d'un torseur sous $\mathfrak{S}_4$, ou
encore d'un rev.\ étale de degré 4 de $S$, ou à celle d'un
\emph{tétraèdre} sur $S$, à celle d'un \emph{octaèdre orienté} sur $S$.
\note{« à celle d'un octaèdre orienté sur $S$ » est écrit à gauche, en bas
de la page, et rattaché à la fin du corollaire}

\page{35}
2) Soit $\Gamma = \mathbb{D}_n$ \emph{($n \geq 3$)}, \struck{Alors}
\add{$\Gamma^{+}$ [$\simeq \mathbb{Z}/n\mathbb{Z}$] le ss-groupe canonique
d'indice 2, $T = \mathfrak{Z}(\Gamma^{+})$ l'unique tore maximal qui}
\[
  \mathfrak{Z}(\Gamma) = {}_2T \simeq \lbrace \pm 1 \rbrace
\]
\[
  \mathcal{N}(\Gamma) =
  \begin{cases}
    \Gamma & \text{si } n \text{ pair} \\
    \Gamma \cdot {}_2T & \text{si } n \text{ impair}
  \end{cases}
\]
\[
  \mathcal{N}(\Gamma)/\mathfrak{Z}(\Gamma) \simeq \Gamma/{}_2\Gamma^{+} \subset \mathrm{Aut}(\Gamma) \,]
\]
\note{devant $\mathfrak{Z}(\Gamma) =$ un premier membre biffé et
noirci ; devant $\mathcal{N}(\Gamma)$, ajouté à gauche,
« $\mathrm{Aut}(G, \Gamma) =$ ». Entre $\mathfrak{Z}(\Gamma)$ et
$\Gamma/{}_2\Gamma^{+}$ de la dernière ligne, une formule biffée
(« $\simeq \mathrm{Aut}(\Gamma) \ldots$ »), et au-dessus de l'inclusion un
signe biffé. En marge gauche, un mot non lu avec une flèche, un
« vérifier » (?) encerclé, et un cercle marqué de points, « $n \geq 3$ »,
non redessiné}

\textbf{\emph{Cor Main}} \add{Sur $S$ sch.\ de car.\ 0, si $n$ pair} La cat.\
des $(G, \Gamma)$, $G$ forme de $\mathrm{GP}(1)$, $\Gamma \subset G$
\struck{sous-groupe fini étale} \add{forme de $\mathbb{D}_n$} ($n$ fixé
$\geq 3$) équivaut à celle des rev.\ principaux de $S$ de groupe
$\mathbb{D}_n$, \struck{\ill{}} \add{(ou encore, \ill{} $\equiv$ \ill{}
racines primitives $n$-ièmes de 1 \ill{})}, i.e.\ celle des
\emph{polygones réguliers} sur $S$ de degré $n$.
\note{« $\geq 3$ » est surchargé, peut-être en « $\geq 4$ »}

Si $n$ est impair $\geq 3$, il faut faire le produit de la cat.\ des rev.\
principaux par celle des rev.\ \ill{} \ill{} \ill{} \ill{}
$\Gamma_S$ \ill{} \ill{} \ill{} que pour $\Gamma' = \mathbb{D}_{2n}$.
\note{ce paragraphe est entouré d'un trait qui le rattache au corollaire}

\marginal{\ill{} \ill{} \ill{} \ill{}\ldots{} \ill{} de l'étude d'une
\ill{} \ill{} \ill{} « \ill{} » \ill{} \ill{}}
\marginal{conséquences \ill{} \ill{} \ill{} d'une \ill{} \ill{}
satisfaisant \ill{} \ill{}\ldots}
\note{à gauche du paragraphe, un griffonnage en hachures serrées couvre
plusieurs mots d'une note marginale oblique, dont on ne lit que des
fragments}

4 \& 5) Soit $\Gamma \simeq \mathfrak{A}_4$ ou $\mathfrak{S}_4$
\[
  \mathfrak{Z}(\Gamma) = 1
\]
\[
  \mathrm{Aut}(G, \Gamma) = \mathcal{N}(\Gamma) = \mathcal{N} \simeq \mathfrak{S}_4
\]
\[
  \mathcal{N}(\Gamma)/\mathfrak{Z}(\Gamma) \simeq \mathcal{N}(\Gamma) \xrightarrow{\sim} \mathrm{Aut}(\Gamma)
\]

\textbf{\emph{Cor.}} 1-1 entre \struck{sous-groupes} formes de
$\mathbb{D}_2$, de $\mathfrak{A}_4$, de $\mathfrak{S}_4$, contenues dans
$G$ fixé (formes \ill{} $=$ espaces homogènes sous $G$)

\textbf{\emph{Cor}} Équivalence de catégories, sur base $S$ de car.\ 0, des
$(G, \Gamma)$, $G$ forme de $\mathrm{GP}(1)$, $\Gamma$ forme de
$\mathbb{D}_2$, ou $\mathfrak{A}_4$, ou $\mathfrak{S}_4$, et avec la cat.\
des torseurs sous $\mathfrak{S}_{4S}$, ou celle des rev.\ étales de degré 4,
ou celle des tétraèdres au-dessus de $S$, ou celle des octaèdres au-dessus
de $S$, ou des formes de $\mathfrak{A}_4$, ou de $\mathfrak{S}_4$, sur $S$.
\note{les deux corollaires sont marqués d'un trait vertical en marge}

\page{36}
\struck{5) Soit $\Gamma = \mathfrak{A}_5$. On a
$\mathfrak{Z}(\Gamma) = \lbrace 1 \rbrace$,
\add{$\mathrm{Aut}(G, \Gamma) =$} $\mathcal{N}(\Gamma) = \Gamma$,
$\mathcal{N}(\Gamma)/\mathfrak{Z}(\Gamma) = \Gamma \hookrightarrow \mathrm{Aut}(\Gamma) \simeq \mathfrak{S}_5$}
\note{le bloc 5) est encadré et barré de trois traits obliques ; il est
transcrit lisiblement}

Dans chacun de ces cas, on a un représentant canonique de la catégorie
fibrée envisagée, défini sur $\mathbb{Q}$ : \struck{\ill{} $\mathbb{Z}$}

\emph{Cas} 1) $\mu_n$ opère sur $\mathbb{G}_{m\mathbb{Q}}$ donc sur
$\widehat{\mathbb{G}}_{m\mathbb{Q}} \simeq \mathbb{P}^1_{\mathbb{Q}}$,
[l'opération s'obtient aussi en faisant opérer $\mu_{2n}$ sur
$\underline{O}^2_{\mathbb{Q}}$, par la façon tautologique sur le premier
facteur $\underline{O}_{\mathbb{Q}}$, et \struck{\ill{}} \add{trivialement}
sur le deuxième, et en passant à
$\mathbb{P}^1_{\mathbb{Q}} = \check{\mathbb{P}}(\underline{O}^2_{\mathbb{Q}})$],
\add{ce qui} fait jouer au quotient de $\mu_{2n}$ : $\mu_n$ \ldots]
\marginal{\textbf{NB} Il n'y a \ill{} de formes non \emph{tordues} de
$\mathbb{Z}/n\mathbb{Z}$ contenues dans une $G$ forme de $\mathrm{GP}(1)$ sur
$\mathbb{Q}$, que si l'équation ½ cyclotomique de degré $n$ \ill{} \ill{}
\struck{\ill{} \ill{} \ill{}} \struck{sur $\mathbb{Q}$ a \ill{}}, ce qui
signifie que $n \in \lbrace 2, 3, 4, 6 \rbrace$. Itou pour avoir une forme
non tordue de $\mathbb{D}_n$ dans une $G$. Pour avoir une forme standard
$\Gamma \subset G$ \ill{} \ill{} \ill{} \ill{} \ill{} !}

\emph{Cas} 2, 3 $\mathbb{Z}/2 \cdot_{1/2} \mu_n$ opère sur
$\mathbb{G}_{m\mathbb{Q}}$ donc sur
\[
  \widehat{\mathbb{G}}_{m\mathbb{Q}} = \mathbb{P}^1_{\mathbb{Q}} .
\]
[\textbf{ReNB} Sur un corps de base (disons) donné $k$, et $n$ donné,
conditions équiv.\ a) $\exists$ $\alpha \in k$ solution de l'équation
« cyclotomique » $\Psi_n(\alpha) = 0$ b) $\exists$ forme non tordue de
$\mathbb{Z}/n\mathbb{Z}$ $\subset$ forme $G$ de $\mathrm{GP}(1)$ (alors $G$
est néc.\ non tordue) c) Itou avec $\mathbb{D}_n$]

\struck{Cas} En fait, tous ces cas sont couverts par les plongements
naturels de $\mu_n$, $\mathbb{Z}/2 \cdot_{1/2} \mu_n$ dans
$\mathbb{Z}/2 \cdot_{1/2} \mathbb{G}_m$, normalisateur du tore maximal
standard de
$\mathrm{GP}(1)_{\mathbb{Q}} \simeq \mathrm{GL}(2)_{\mathbb{Q}}/\mathbb{G}_{m\mathbb{Q}}$.

\emph{Cas} 4, 5 \add{prenons} On prend le plongement principal
$\mathbb{D}_{2\mathbb{Q}} = \mathbb{Z}/2 \cdot_{1/2} \mu_{2\mathbb{Q}}$ de
$\mathrm{GP}(1)$, on prend le \uncertain{normalisateur} de
$\mathbb{D}_{2\mathbb{Q}}$ dans $\mathrm{GP}(1)$ --- $\mathfrak{S}_4$
\ill{} : fait une forme \add{(tordue ?)} de $\mathfrak{S}_4$ plongée dans
$\mathrm{GP}(1)$, il faudrait regarder de plus près. [\ill{} \ill{}
\ill{} \ill{} \ill{} une forme de $\mathfrak{A}_4$,
\marginal{à chercher \ill{} un autre représentant plus commode ? \ill{}
\ill{} \ill{} \ill{}}
\note{les cas 4, 5 sont marqués d'un trait vertical en marge ; la note
marginale est en partie couverte par un griffonnage}

\page{37}
6) $\Gamma \simeq \mathfrak{A}_5$.
\[
  \mathfrak{Z}(\Gamma) = \lbrace 1 \rbrace, \qquad
  \mathcal{N}(\Gamma) = \Gamma \simeq \mathfrak{A}_5
\]
\emph{Une fois} choisi un cas de référence \add{$(G_0, \Gamma_0)$} sur
$\mathbb{Q}$, si $S$ est un schéma de car.\ 0, se donner un couple
$(G \supset \Gamma)$, avec $G$ forme de $\mathrm{GP}(1)$, $\Gamma$ forme de
$\mathfrak{A}_5$, revient au même que de se donner une forme
\emph{intérieure} de $\Gamma_0$, ou un torseur sous $\Gamma_0$. Si on
\add{choisit} \struck{\ill{}} $\Gamma_{0k} \simeq \mathfrak{A}_5$, il
revient au même de se donner \struck{\ill{}} un rev.\ de degré 5 de $S$ +
trivialisation du rev.\ quadratique de ses orientations.

Mais on a ceci : $\exists$ \add{sur $k$} forme non tordue de
$\mathfrak{A}_5$ dans une $G$ (forme de $\mathrm{GP}(1)$) sss l'équation
½ cyclotomique
\[
  \alpha^2 + \alpha - 1 = 0
\]
a ses racines dans $k$ \struck{et alors s'il existe une fo} \add{Pour que}
de plus $G$ soit non tordue, il f.\ et s.\ que $k$ contienne les racines
cinquièmes de 1.
\marginal{vérifier}
\note{« vérifier » est souligné deux fois ; la phrase « Pour que \ldots{}
de 1 » est marquée d'un trait vertical}

Sur un $k$ \emph{muni} d'une solution $\alpha$, il y a un plongement
\struck{\ill{}} de $\mathfrak{A}_5$ dans une $G$ [défini à iso unique près
\add{en exigeant que l'image d'un $5$-cycle permutation circulaire standard
soit $\alpha$ \uncertain{certain} $\alpha$ !]} --- OPS
$k = k_0 = \mathbb{Q}[T]/(T^2 + T - 1) \simeq \mathbb{Q}[\alpha_0]$, où
\[
  \alpha_0 = \tfrac{1}{2}(-1 + \sqrt{5}) \,]
\]
Sur les $k_0$-schémas, la donnée d'un $G \supset \Gamma$, $\Gamma$
\add{$(G_1)$} forme de $\mathfrak{A}_5$ ($\mathrm{GP}(1)$) équivaut à la
donnée d'un torseur sous $\mathfrak{A}_5$, i.e.\ \ldots
\note{le « $+$ » de $\alpha_0$ est entouré. « $(G_1)$ » au-dessus de
$\Gamma$ est une lecture douteuse. La page s'arrête sur « i.e. » et un
tiret ; le bas du feuillet est blanc}

\page{38}
\subsection*{Formes de $\mathbb{D}_n$}
\note{titre souligné, de sa main, en tête de la page 38}

Soit $\Gamma$ un groupe, \struck{tel que} satisfaisant les conditions
suivantes : il existe un ss-groupe $\Gamma^{+}$ de $\Gamma$ tel que
\begin{enumerate}
  \item[a)] $\Gamma^{+}$ d'indice 2 dans $\Gamma$, i.e.\ $\Gamma$ ext.\ de
  $\mathbb{Z}/2 \simeq \Gamma/\Gamma^{+}$ par $\Gamma^{+}$
  \item[b)] $\Gamma^{+}$ \struck{\ill{}} commut.\ [de
  $\Gamma/\Gamma^{+} = \mathbb{Z}/2$ opérant sur $\Gamma^{+}$]
  \item[c)] $\mathbb{Z}/2$ opère sur $\Gamma^{+}$ par $x \mapsto -x$
  \item[e)] $\Gamma^{+} \setminus {}_2(\Gamma^{+})$ engendre le groupe
  $\Gamma^{+}$. [\textbf{NB} si $\Gamma^{+}$ est cyclique, \struck{ce}
  d'ordre $n$, ceci signifie $n \neq 1$]
  \item[d)] $\exists$ $\sigma \in {}_2\Gamma$, $\sigma \notin \Gamma^{+}$
  [i.e.\ l'extension de $\mathbb{Z}/2$ par $\Gamma^{+}$ est scindée].
\end{enumerate}
\note{la lettre e) surcharge un premier signe ; une flèche courbe la
relie à d), sans qu'on voie si elle doit passer après. ${}_2\Gamma$ note
l'ensemble des éléments d'ordre divisant 2}

\struck{On a alors}

\struck{Alors le} Soit \struck{Alors} $\Gamma^{-} = \Gamma \setminus
\Gamma^{+}$. On voit tout de suite \struck{que} (utilisant c) \& d))
\[
  \Gamma^{-} \subset {}_2\Gamma
\]
donc
\[
  {}_2\Gamma = \Gamma^{-} \amalg {}_2\Gamma^{+}
\]
donc
\[
  \Gamma \setminus {}_2\Gamma = \Gamma^{+} \setminus {}_2\Gamma^{+}
\]
donc, utilisant e)
\[
  (*) \qquad \Gamma^{+} = \text{ss-grpe de } \Gamma \text{ engendré par } \Gamma \setminus {}_2\Gamma
\]
ce qui montre donc que $\Gamma^{+}$ est déterminé de façon unique.
D'ailleurs, $\Gamma^{-}$ est un torseur à gauche sous $\Gamma^{+}$, et la
multiplication sur $\Gamma = \Gamma^{+} \amalg \Gamma^{-}$ se déduit de la
structure de groupe sur $\Gamma^{+}$ et de la structure de
$\Gamma^{+}$ tors.\ à g.\ sur $\Gamma^{-}$, grâce :
\begin{enumerate}
  \item[1°)] La multiplication de $\Gamma$ dans
  $\Gamma^{+} \times \Gamma^{+}$ est celle de $\Gamma^{+}$
  \item[2°)] --- \quad $\Gamma^{-} \times \Gamma^{-}$ est l'application
  ct.\ $\to 1 \in \Gamma^{+}$
  \item[3°)] --- \quad $\Gamma^{+} \times \Gamma^{-}$ est induite par
  l'opération de $\Gamma^{+}$ sur $\Gamma^{-}$
  \item[4°)] --- \quad $\Gamma^{-} \times \Gamma^{+}$ est donnée par
  $(\sigma, v) \mapsto \sigma v = (-v) \cdot \sigma$
\end{enumerate}
\note{les tirets rendent ses traits de répétition pour « La multiplication
de $\Gamma$ dans »}

\page{39}
Inversement, partant d'un groupe abélien $\Gamma^{+}$, et d'un torseur
$\Gamma^{-}$ sous $\Gamma^{+}$, on définit sur
$\Gamma = \Gamma^{+} \amalg \Gamma^{-}$ une \struck{structure} loi de
composition par les formules précédentes, et on voit de suite que celle-ci
est une loi de groupe, qui fait de $\Gamma$ un groupe admettant $\Gamma^{+}$
comme sous-groupe \add{abélien} d'indice 2, l'extension de $\mathbb{Z}/2$
par $\Gamma^{+}$ ainsi obtenue donnant lieu : \add{l'opération
$x \mapsto -x$} \struck{\ill{} \ill{}} de $\Gamma^{+}$, et étant scindée
--- à \uncertain{n'importe quel} élément de $\Gamma^{-}$. [Donc $\Gamma$,
$\Gamma^{+}$ satisfont les conditions a) à d) ci-dessus, mais pas néc.\ e)]
On trouve donc un foncteur qui va de la cat.\ des torseurs sous groupes
abéliens (avec les morphismes habituels, par ex.\ iso) vers la catégorie
des ext.\ « diédrales » de $\mathbb{Z}/2$ par groupes abéliens --- et c'est
une équiv.\ de catégories --- et la restriction de ce foncteur \add{à la
ss-cat.\ pleine des} \struck{aux} couples $(\Gamma^{+}, \Gamma^{-})$, avec
$\Gamma^{+}$ satisfaisant e) ci-dessus, \add{et les seuls iso)} est un
foncteur pl.\ fidèle vers la cat.\ des groupes $\Gamma$ satisfaisant
\struck{les} conditions du début, avec les iso. En particulier

\textbf{\emph{Prop}} Soit $n \geq 3$ un entier. La \struck{cat.}\
\add{\uncertain{\emph{gerbe}}} \struck{des} des formes de $\mathbb{D}_n$ est
équivalente à la \struck{catégorie} \add{\uncertain{gerbe}} des droites
affines sur l'anneau $\mathbb{Z}/n$. En particulier, on trouve
\[
  \mathrm{Aut}(\mathbb{D}_n) \simeq \mathrm{Aff}(1, \mathbb{Z}/n) \simeq (\mathbb{Z}/n)^{*} \cdot_{1/2} \mathbb{Z}/n
\]
\note{la proposition est marquée d'un trait vertical en marge. Le mot
ajouté deux fois au-dessus de « cat. » biffé, lu « gerbe », est incertain}

\page{40}
\textbf{\emph{Théorème}} Soient $S$ un schéma. Pour toute forme $G$ de
$\mathrm{GP}(1)_S$ et tout \emph{mono}
\[
  \varphi : \mathfrak{A}_{5S} \hookrightarrow G
\]
soit
\[
  \alpha(\varphi) = A(\varphi(\pi))
\]
où
\[
  A : G \longrightarrow \mathbb{E}^1_S
\]
est l'hom.\ invariant et où $\pi$ est la permutation circulaire standard de
$\mathfrak{A}_5$. Ceci posé : \struck{la cat.\ des} \struck{foncteur}
\struck{couples} \add{de la catégorie-groupoïde des $(G, \varphi)$}
\[
  (G, \varphi) \longmapsto \alpha(\varphi)
\]
(où $G$ est \struck{\ill{} iso} une forme de $\mathrm{GP}(1)_S$ et
$\varphi : \mathfrak{A}_{5S} \hookrightarrow G$, \struck{\ill{}
\emph{discrète} (i.e.\ \ill{} \ill{} autom.\ et \ill{} identité)}),
\struck{et l'application $(G, \varphi) \mapsto \alpha(\varphi)$ définit une
bijection de l'ens.\ des \ill{} classes d'iso, sur l'ens \struck{des}
$\mathrm{Hom}(S, M)$ des $\alpha \in \Gamma(S, \underline{O}_S)$
satisfaisant $\alpha^2 + \alpha - 1 = 0$} \add{dans la catégorie}
discrète définie par l'ens des $\alpha \in \Gamma(S, \underline{O}_S)$ tels
que
\[
  \alpha^2 + \alpha - 1 = 0 ,
\]
est une équivalence de catégories. En d'autres termes : la catégorie des
$(G, \varphi)$ est \emph{rigide}, et pour $S$ variable, le foncteur
$S \mapsto F(S) =$ \struck{ens.\ des} classes d'iso de couples
$(G, \varphi)$ sur $S$, est représentable par
\[
  M = \mathrm{Spec}\, \mathbb{Z}[T]/(T^2 + T - 1) .
\]
grâce à $\varphi \mapsto \alpha(\varphi)$.
\note{le passage biffé est encadré et barré de traits obliques ; le texte
qui le remplace (« dans la catégorie discrète \ldots ») est écrit sous le
cadre}
\marginal{il suffit que \struck{\ill{}} \uncertain{soit trivial} (\ill{}
\ill{} \ill{}) sur les fibres\ldots}
\marginal{On pose $\alpha'(\varphi) = -1 - \alpha(\varphi) =
-1/\alpha(\varphi)$}
\note{la première note marginale est en face de l'énoncé, la seconde en
face de « Ceci posé »}

\textbf{\emph{Cor}} Soit $(G, \varphi)$ comme dessus, et soit $u$ un
automorphisme de $\mathfrak{A}_5$. \struck{Pour qu'on ait} Si $u$ est
intérieur on a
$\alpha(\varphi \circ u) = \alpha(\varphi)$ i.e.\
$(G, \varphi) \sim (G, \varphi \circ u)$. Si $u$ n'est intérieur sur aucune
fibre, on a $\alpha(\varphi \circ u) = \alpha'(\varphi)$
\note{fin du lot ; la page s'arrête sans ponctuation finale}

\end{document}
